\documentclass[12pt,a4paper]{article} % Definimos tamanho da letra e papel
\usepackage[utf8]{inputenc}
\usepackage[brazilian]{babel}
\usepackage[T1]{fontenc}
\usepackage{graphicx} 
\usepackage{xurl} % Para os links do GitHub não saírem da margem
\usepackage[hidelinks]{hyperref} % Para links clicáveis
\usepackage[margin=2.5cm]{geometry} % Ajusta as margens para padrão acadêmico
\usepackage[table,xcdraw]{xcolor}
\usepackage{array}
\usepackage{booktabs}
\usepackage{longtable}
\usepackage{seqsplit}




\title{Adverse drug reaction records from the Brazilian pharmacovigilance system (VigiMed), 2018--2025}

\author{
    Ana Carolina Jacoby$^{1, *}$, 
    Silmara Gomes Barnabé$^2$, 
    Mell Amisa Matsuda$^1$, \\
    Mariana Recamonde Mendoza$^2$, 
    Sílvio César Cazella$^1$
}

\date{\vspace{-5ex}} % Remove a data automática para ganhar espaço

\begin{document}

\maketitle

% --- Afiliações e Contato ---
\begin{center}
    \footnotesize
    $^1$ Federal University of Health Sciences of Porto Alegre (UFCSPA) \\
    $^2$ Federal University of Rio Grande do Sul (UFRGS) \\
    $^*$ Corresponding author: \texttt{ana.jacoby@ufcspa.edu.br}
\end{center}

\begin{abstract}

A farmacovigilância é fundamental para o monitoramento contínuo da segurança de produtos farmacêuticos. No Brasil, o sistema VigiMed centraliza as notificações nacionais de eventos adversos. Este trabalho apresenta um conjunto de dados unificado com informações clínicas, farmacológicas e sociodemográficas de suspeitas de reações adversas a medicamentos (RAMs) e eventos adversos pós-vacinação (EAPV). Os dados derivam de relatos espontâneos submetidos à Agência Nacional de Vigilância Sanitária (ANVISA) entre 2018 a 2025. O conjunto original compreendia [X] registros e [Y] atributos. Após as etapas de pré-processamento e limpeza, foi obtido um dataset final com [Z] registros e [W] atributos. Os dados disponibilizados visam apoiar a comunidade científica em estudos de segurança de medicamentos, incluindo a identificação de potenciais sinais de segurança e a análise de desfechos clínicos. Além disso, o conjunto pode ser reutilizado no desenvolvimento e validação de abordagens metodológicas, como modelos de aprendizado de máquina, aplicados à avaliação de gravidade e à vigilância em saúde pública.
\end{abstract}


\section{Background & Summary}
Seu texto aqui.

\section{Methods}

Os dados foram coletados do sistema nacional brasileiro de farmacovigilância, VigiMed, mantido pela Agência Nacional de Vigilância Sanitária (ANVISA). O VigiMed é o sistema oficial para a coleta de relatos espontâneos de suspeitas de reações adversas a medicamentos (RAMs) e vacinas, submetidos por profissionais de saúde, pacientes e empresas farmacêuticas em todo o território nacional.
O conjunto de dados compreende relatos espontâneos relacionados a medicamentos e vacinas notificados entre 2018 e 2025. Os dados brutos foram obtidos a partir do Portal de Dados Abertos da ANVISA \url{https://dados.anvisa.gov.br/dados/} e posteriormente processados para a construção do conjunto de dados utilizado neste estudo.

\subsection{Ethical Approval} 

De acordo com a Resolução nº 510/2016 do Conselho Nacional de Saúde, este estudo está dispensado de aprovação por Comitê de Ética em Pesquisa, uma vez que utiliza dados secundários de acesso público, agregados e anonimizados. A estrutura do conjunto de dados inclui informações sobre eventos adversos, produtos suspeitos e concomitantes, cronologia das notificações, critérios de gravidade e variáveis sociodemográficas selecionadas. Em conformidade com a Lei Geral de Proteção de Dados (LGPD), não são disponibilizados identificadores diretos ou informações que permitam a identificação dos indivíduos.

\section{Data Records}

O conjunto de dados processado está disponível no repositório Dataverse \textbf{[inserir link]}, sob o DOI \textbf{[inserir DOI].} Os dados são fornecidos em formato CSV e abrangem notificações espontâneas de reações adversas a medicamentos e vacinas no Brasil entre 2018 e 2025. O dataset está organizado em uma arquitetura em estrela, composta por três tabelas de fatos e tabelas de dimensões compartilhadas, permitindo análises em diferentes níveis de granularidade.

Fato notificações: Esta tabela contém informações detalhadas de cada notificação, incluindo a identificação da notificação, relação entre medicamentos e eventos adversos, detentor do registro, componente suspeito, ação adotada e cronologia da notificação. Também inclui indicações médicas, dosagem, duração do tratamento, início e fim da administração, forma farmacêutica, via de administração e número do lote. Chaves primárias conectam cada notificação aos registros de medicamentos e aos códigos ATC.

Fato medicamentos: Fornece informações detalhadas sobre os medicamentos e vacinas notificadas em cada registro, incluindo nome do produto, princípios  ativos, detentor do registro,  código ATC, concentração, dose, forma e via de  administração, além de indicações  médicas  e problemas relacionados ao medicamento.

Fato reações: Registra as reações adversas  associadas a cada  notificação, incluindo  datas  de início e término, duração,  gravidade, desfecho e indicadores de eventos graves, como óbito,  ameaça a  vida e hospitalização. Os eventos são codificados seguindo a terminologia MedDRA, permitindo a análise em vários níveis desta hierarquia.

As tabelas de dimensões compartilhadas contextualizam as características sociodemográficas dos pacientes, além de atributos temporais e geográficos. Os medicamentos e as reações adversas seguem as classificações ATC/DDD e a terminologia MedDRA, respectivamente. Contudo, devido a restrições de licenciamento, os dicionários originais não são redistribuídos.

Essa organização permite análises integradas, desde o nível de notificação individual até padrões de uso de medicamentos e ocorrência de eventos adversos, garantindo flexibilidade para diferentes abordagens analíticas.

\textbf{SUGIRO TABELAS  FATOS DESCRITAS
EX:}


\begin{longtable}{
|>{\raggedright\arraybackslash}p{5.5cm}
|>{\raggedright\arraybackslash}p{2.5cm}
|>{\raggedright\arraybackslash}p{6.5cm}|}

\caption{Tabela Fato Reações (fat\_reacoes)} \\
\hline
\textbf{Variable} & \textbf{Data type} & \textbf{Description} \\
\hline
\endfirsthead

\hline
\textbf{Variable} & \textbf{Data type} & \textbf{Description} \\
\hline
\endhead

\seqsplit{IDENTIFICACAO\_NOTIFICACAO} & object &
Unique identifier for the adverse event notification (e.g., BR-ANVISA-300000004). \\ \hline

\seqsplit{HASH\_BRONZE} & object &
SHA-256 hash of the original raw data columns for change detection. \\ \hline

\seqsplit{HASH\_SILVER} & object &
SHA-256 hash of the transformed data columns for change detection. \\ \hline

\seqsplit{DATA\_INICIO\_HORA} & datetime64 &
Normalized date and time when the adverse reaction started. \\ \hline

\seqsplit{DATA\_FINAL\_HORA} & datetime64 &
Normalized date and time when the adverse reaction ended or resolved. \\ \hline

\seqsplit{ATUALIZADO} & object &
Date when the notification record was last updated. \\ \hline

\seqsplit{LLT\_CHAVE} & int64 &
Foreign key to the MedDRA Low Level Term (LLT) in dim\_soc\_llt dimension table. \\ \hline

\seqsplit{GRAVE\_CHAVE} & Int64 &
Severity indicator:
0 = Unknown;
1 = Yes (Serious);
2 = No (Not Serious). \\ \hline

\seqsplit{DESFECHO\_CHAVE} & Int64 &
Outcome code:
0 = Unknown;
1 = Not Recovered/Unresolved/Ongoing;
2 = Recovering/Resolving;
3 = Recovered/Resolved;
4 = Fatal/Death;
5 = Recovered/Resolved with sequelae. \\ \hline

\seqsplit{DURACAO\_TIPO\_CHAVE} & int64 &
Duration unit code:
0 = Unknown;
1 = Second(s);
2 = Minute(s);
3 = Hour(s);
4 = Day(s);
5 = Week(s);
6 = Month(s);
7 = Year(s);
8 = Decade(s);
9 = Cycle(s). \\ \hline

\seqsplit{DURACAO\_VALOR} & float64 &
Numerical value of the adverse reaction duration. Interpret together with DURACAO\_TIPO\_CHAVE. \\ \hline

\seqsplit{GRAVIDADE\_RESULTADO\_OBITO} & int64 &
Reaction resulted in death:
1 = Yes; 0 = No. \\ \hline

\seqsplit{GRAVIDADE\_AMEACA\_VIDA} & int64 &
Life-threatening reaction:
1 = Yes; 0 = No. \\ \hline

\seqsplit{GRAVIDADE\_INCAPACIDADE} & int64 &
Persistent or significant incapacity/disability:
1 = Yes; 0 = No. \\ \hline

\seqsplit{GRAVIDADE\_HOSPITALIZACAO} & int64 &
Required hospitalization or prolonged hospitalization:
1 = Yes; 0 = No. \\ \hline

\seqsplit{GRAVIDADE\_OUTRO\_EFEITO} & int64 &
Other clinically significant effect:
1 = Yes; 0 = No. \\ \hline

\seqsplit{GRAVIDADE\_ANOMALIA\_CONGENITA} & int64 &
Congenital anomaly or birth defect:
1 = Yes; 0 = No. \\ \hline

\end{longtable}



\begin{longtable}{
|>{\raggedright\arraybackslash}p{5.5cm}
|>{\raggedright\arraybackslash}p{2.5cm}
|>{\raggedright\arraybackslash}p{6.5cm}|}

\caption{Tabela Fato Medicamentos (fat\_medicamentos)} \\
\hline
\textbf{Variable} & \textbf{Data type} & \textbf{Description} \\
\hline
\endfirsthead

\hline
\textbf{Variable} & \textbf{Data type} & \textbf{Description} \\
\hline
\endhead

\seqsplit{IDENTIFICACAO\_NOTIFICACAO} & object &
Unique identifier for the adverse event notification. \\ \hline

\seqsplit{NOME\_MEDICAMENTO\_WHODRUG} & object &
Drug name according to WHODrug dictionary. \\ \hline

\seqsplit{PRINCIPIOS\_ATIVOS\_WHODRUG} & object &
Active ingredients according to WHODrug dictionary. \\ \hline

\seqsplit{ATC\_CODE\_LEVEL\_4} & object &
ATC (Anatomical Therapeutic Chemical) code at level 4. Disaggregated from pipe-delimited values. \\ \hline

\seqsplit{DETENTOR\_REGISTRO\_SCORE} & float64 &
Fuzzy matching score (0-100) for marketing authorization holder. \\ \hline

\seqsplit{DETENTOR\_REGISTRO\_CHAVE} & Int64 &
Foreign key to marketing authorization holder dimension table. 0 = Unknown. \\ \hline

\seqsplit{CONCENTRACAO\_TIPO\_CHAVE} & Int64 &
Concentration unit code (e.g., mg/ml, \%, UI/ml). \\ \hline

\seqsplit{CONCENTRACAO\_TIPO\_VALOR} & object &
Concentration unit label (e.g., "mg/ml", "percentage"). \\ \hline

\seqsplit{CONCENTRACAO\_VALOR} & float64 &
Calculated concentration value (numerator/denominator ratio). \\ \hline

\seqsplit{CONCENTRACAO\_VALOR\_NUMERADOR} & float64 &
Numerator value of concentration (e.g., 500 in "500mg/ml"). \\ \hline

\seqsplit{CONCENTRACAO\_VALOR\_DENOMINADOR} & float64 &
Denominator value of concentration (e.g., 1 in "500mg/1ml"). \\ \hline

\seqsplit{PROB\_ADIC\_ERRO\_DE\_MEDICACAO} & int64 &
Medication error: 1 = Yes; 0 = No. \\ \hline

\seqsplit{PROB\_ADIC\_EXPOSICAO\_OCUPACIONAL} & int64 &
Occupational exposure: 1 = Yes; 0 = No. \\ \hline

\seqsplit{PROB\_ADIC\_FALSIFICACAO} & int64 &
Counterfeit/falsification: 1 = Yes; 0 = No. \\ \hline

\seqsplit{PROB\_ADIC\_LOTES\_TESTADOS\_E\_DENTRO\_DAS\_ESPECIFICACOES} & int64 &
Tested batches within specifications: 1 = Yes; 0 = No. \\ \hline

\seqsplit{PROB\_ADIC\_LOTES\_TESTADOS\_E\_FORA\_DAS\_ESPECIFICACOES} & int64 &
Tested batches out of specifications: 1 = Yes; 0 = No. \\ \hline

\seqsplit{PROB\_ADIC\_MEDICAMENTO\_TOMADO\_FORA\_DA\_DATA\_DE\_VALIDADE} & int64 &
Drug taken after expiration date: 1 = Yes; 0 = No. \\ \hline

\seqsplit{PROB\_ADIC\_MEDICAMENTO\_TOMADO\_PELO\_PAI} & int64 &
Drug taken by the father: 1 = Yes; 0 = No. \\ \hline

\seqsplit{PROB\_ADIC\_SUPERDOSE} & int64 &
Overdose: 1 = Yes; 0 = No. \\ \hline

\seqsplit{PROB\_ADIC\_USO\_INCORRETO} & int64 &
Incorrect use: 1 = Yes; 0 = No. \\ \hline

\seqsplit{PROB\_ADIC\_USO\_OFF\_LABEL\_USO\_SEM\_REGISTRO} & int64 &
Off-label use or unregistered use: 1 = Yes; 0 = No. \\ \hline

\seqsplit{PROB\_ADIC\_ABUSO} & int64 &
Drug abuse: 1 = Yes; 0 = No. \\ \hline

\seqsplit{PK\_LLT\_INDICACAO\_MEDDRA} & Int64 &
Foreign key to MedDRA LLT for reported MedDRA indication (fuzzy matched). \\ \hline

\seqsplit{PK\_LLT\_INDICACAO\_RELATADA\_NOTIFICADOR\_INICIAL} & Int64 &
Foreign key to MedDRA LLT for initial notifier reported indication (fuzzy matched). \\ \hline

\seqsplit{DOSE\_TIPO\_CHAVE} & object &
Dose unit label (e.g., "milligram (mg)", "microgram (mcg)"). \\ \hline

\seqsplit{DOSE\_TIPO\_VALOR} & Int64 &
Dose unit code. 0 = Unknown. \\ \hline

\seqsplit{DOSE\_VALOR} & float64 &
Numerical dose value. \\ \hline

\seqsplit{FREQUENCIA\_DOSE} & object &
Dose frequency description. \\ \hline

\seqsplit{POSOLOGIA} & object &
Posology/dosage regimen description. \\ \hline

\seqsplit{DURACAO\_TIPO\_CHAVE} & int64 &
Duration unit code:
0 = Unknown; 1 = Second(s); 2 = Minute(s); 3 = Hour(s); 4 = Day(s); 5 = Week(s); 6 = Month(s); 7 = Year(s); 8 = Decade(s); 9 = Cycle(s). \\ \hline

\seqsplit{DURACAO\_TIPO\_VALOR} & object &
Duration unit label (e.g., "DIA", "HORA", "SEMANA"). \\ \hline

\seqsplit{DURACAO\_VALOR} & float64 &
Numerical duration value. Interpret with DURACAO\_TIPO\_CHAVE. \\ \hline

\seqsplit{INICIO\_ADMINISTRACAO} & datetime64 &
Normalized date when drug administration started. \\ \hline

\seqsplit{FIM\_ADMINISTRACAO} & datetime64 &
Normalized date when drug administration ended. \\ \hline

\seqsplit{FORMA\_FARMACEUTICA\_SCORE} & float64 &
Fuzzy matching score (0-100) for pharmaceutical form. \\ \hline

\seqsplit{FORMA\_FARMACEUTICA\_CHAVE} & Int64 &
Foreign key to pharmaceutical form dimension table. 0 = Unknown. \\ \hline

\seqsplit{VIA\_ADMINISTRACAO} & object &
Route of administration (original text). \\ \hline

\seqsplit{VIA\_ADMINISTRACAO\_CHAVE} & int64 &
Normalized route of administration code (fuzzy matched). \\ \hline

\seqsplit{VIA\_ADMINISTRACAO\_MAE\_PAI} & object &
Route of administration for mother/father (original text). \\ \hline

\seqsplit{VIA\_ADMINISTRACAO\_MAE\_PAI\_CHAVE} & int64 &
Normalized route code for mother/father (fuzzy matched). \\ \hline

\seqsplit{NUMELO\_LOTE} & object &
Batch/lot number of the drug. \\ \hline

\seqsplit{HASH\_BRONZE} & object &
SHA-256 hash of original raw data columns for change detection. \\ \hline

\seqsplit{RELACAO\_MEDICAMENTO\_EVENTO\_VALOR} & object &
Drug-event relationship description. \\ \hline

\seqsplit{RELACAO\_MEDICAMENTO\_EVENTO\_CHAVE} & Int64 &
Drug-event relationship code:
0 = Unknown; 1 = Suspect; 2 = Concomitant; 3 = Drug not administered; 4 = Interaction. \\ \hline

\seqsplit{COMPONENTE\_SUSPEITO\_VALOR} & object &
Suspected component description. \\ \hline

\seqsplit{COMPONENTE\_SUSPEITO\_CHAVE} & Int64 &
Suspected component code:
0 = Unknown; 1 = Active ingredient; 2 = Excipient; 3 = Solvent; 4 = Dye; 5 = Preservative; 6 = Flavoring agent; 7 = Overage; 8 = Antioxidant; 9 = Stabilizer. \\ \hline

\seqsplit{ACAO\_ADOTADA\_VALOR} & object &
Action taken description. \\ \hline

\seqsplit{ACAO\_ADOTADA\_CHAVE} & Int64 &
Action taken code:
0 = Unknown; 1 = Drug withdrawn; 2 = No dose change; 3 = Not applicable; 4 = Dose reduced; 5 = Dose increased. \\ \hline

\end{longtable}





\section{Data Overview (optional)}

Creio ser importante colocar uma  visão básica dos dados, é  opcional,  então aguardo orientação.

\section{Technical Validation}

\section{Usage Notes (optional)}

\section{Data Availability}

\section*{Code Availability}

The code used for the pre-processing of the dataset presented in this paper is available at: \url{https://github.com/mlab-inf-ufrgs/vigimed}

\section*{Author Contributions}

A.C.J. e S.G.B. contribuíram igualmente para este trabalho. A.C.J.: concepção do estudo, modelagem de dados, curadoria de dados, análise de dados, redação. S.G.B.: concepção do estudo, modelagem de dados, curadoria de dados, análise de dados, redação. M.A.M.: processamento de dados, desenvolvimento de código, validação de dados. M.R.M.: supervisão, orientação metodológica, revisão do manuscrito. S.C.C.: supervisão, orientação metodológica, revisão do manuscrito.

Todos os autores leram e aprovaram a versão final do manuscrito.

\section*{Competing Interests}
The authors declare no competing interests.

\section*{Funding}

This research did not receive specific funding.

\section*{Acknowledgements}
Os autores agradecem ao Centro Colaborador da OMS para Metodologia de Estatísticas de Medicamentos (Oslo, Noruega) pelo uso do sistema ATC/DDD e à equipe do MedDRA pelo acesso à sua terminologia sob licença acadêmica. O uso desses sistemas destina-se exclusivamente à classificação e codificação. As análises e conclusões aqui apresentadas são de inteira responsabilidade dos autores e não implicam endosso por parte dos licenciadores ou das organizações mencionadas.

% --- Referências em Vancouver ---
\bibliographystyle{vancouver}
\bibliography{references} % Procure ter um arquivo chamado referencias.bib

\end{document}